\documentclass{article}
\usepackage{amsmath,amsthm,amssymb}
\usepackage[margin=1in]{geometry}

\DeclareMathOperator{\dinv}{dinv}
\DeclareMathOperator{\di}{di}
\DeclareMathOperator{\RR}{RR}
\DeclareMathOperator{\SSYT}{SSYT}
\newcommand{\DS}{\operatorname{DS}}
\newcommand{\DSstar}{\operatorname{DS}^{\ast}}
\newcommand{\wt}{\operatorname{wt}}
\newcommand{\Zge}{\mathbb Z_{\ge 0}}

\newtheorem{theorem}{Theorem}
\newtheorem{conjecture}{Conjecture}
\newtheorem{remark}{Remark}

\title{Dyck Symmetric Functions}
\author{}
\date{}

\begin{document}
\maketitle

\section{Overview}

This note explains Dyck and dual Dyck symmetric functions and the
\(r=s\tau+1\) rational analogue.  The guiding picture has three levels.
At \(\tau=0\), the theory is a version of the classical RSK and dual RSK
correspondences.  At \(\tau=1\), the ordinary Dyck and dual Dyck symmetric
functions are theorem-level objects: the 2026 paper proves the dual formula by
an explicit tableau insertion bijection and derives the affine formula by the
standard involution on symmetric functions.  For \(\tau>1\), the same
Schur-expansion statements are conjectural.  They are supported here by finite
computations, but no insertion algorithm is currently known for the rational
case.

The terminology ``affine'' and ``dual'' refers to the two local inequalities
allowed inside each factor of a factorization.  The affine inequality is
\(x_{i+1}\le x_i+\tau\).  The dual inequality is
\(x_{i+1}>x_i+\tau\).  When \(\tau=1\), the dual condition becomes
\(x_{i+1}\ge x_i+2\), which is the condition used in the 2026 insertion
algorithm.  When \(\tau=0\), the dual condition becomes strict increase, so
the dual side reduces to the familiar dual-RSK setting.

\section{Step-$\tau$ dinv and rational Dyck factors}

Fix a nonnegative integer \(\tau\).  For a finite integer sequence
\(w=(w_0,\ldots,w_{\ell-1})\), define
\[
  \dinv_\tau(w)=\sum_{0\le i<j<\ell} d_\tau(w_i,w_j),
\]
where
\[
  d_\tau(a,b)=
  \begin{cases}
    \max(0,a+\tau-b),& a\le b,\\
    \max(0,b+1+\tau-a),& a>b.
  \end{cases}
\]
The rational family considered here corresponds to slopes
\[
  r=s\tau+1.
\]
This note does not attempt to treat arbitrary coprime rational slopes.

A finite integer sequence is called an \emph{affine \(\tau\)-Dyck sequence} if
\[
  w_{i+1}\le w_i+\tau
  \qquad(0\le i<\ell-1),
\]
and a \emph{dual \(\tau\)-Dyck sequence} if
\[
  w_{i+1}>w_i+\tau
  \qquad(0\le i<\ell-1).
\]
The empty sequence and one-entry sequences satisfy both conditions vacuously.

Let \(S\) be a finite multiset of integers.  A factorization of \(S\) is a
sequence
\[
  \mathcal F=(F_0,F_1,F_2,\ldots)
\]
of finite words, all but finitely many empty, whose concatenation is a
rearrangement of \(S\).  It is affine, respectively dual, if every factor is an
affine, respectively dual, \(\tau\)-Dyck sequence.  We use zero-indexed
symmetric-function variables
\(\mathbf x=(x_0,x_1,x_2,\ldots)\), and write
\[
  x^{\mathcal F}=\prod_{i\ge0} x_i^{|F_i|}.
\]
The affine and dual rational Dyck symmetric functions are
\[
  \DS^{(\tau)}(S,d;\mathbf x)
  =
  \sum_{\substack{\mathcal F\text{ affine }\tau\text{-Dyck factorization of }S\\
                  \dinv_\tau(F_0F_1F_2\cdots)=d}}
       x^{\mathcal F}
\]
and
\[
  \operatorname{DS}^{\ast,(\tau)}(S,d;\mathbf x)
  =
  \sum_{\substack{\mathcal F\text{ dual }\tau\text{-Dyck factorization of }S\\
                  \dinv_\tau(F_0F_1F_2\cdots)=d}}
       x^{\mathcal F}.
\]

\section{Rational Dyck tableaux}

A rational Dyck tableau of step \(\tau\) is a left-aligned tableau whose row
lengths form a partition.  Rows are indexed from top to bottom.  Each row,
read left to right, is required to be a dual \(\tau\)-Dyck sequence.  Columns
are read from bottom to top and are required to satisfy the affine
\(\tau\)-Dyck inequality.  Equivalently, if row \(i\) is immediately above row
\(i+1\), then entries in a common column satisfy
\[
  P_i[j]\le P_{i+1}[j]+\tau.
\]
The row-reading word is read from the bottom row upward, each row left to
right:
\[
  \RR(P)=P_{\mathrm{bottom}}P_{\mathrm{next}}\cdots P_{\mathrm{top}}.
\]
The shape \(\lambda(P)\) is the partition of row lengths.

For \(\tau=1\), this is the Dyck tableau convention used in the classical
paper.  For \(\tau>1\), the same definition supplies the conjectural rational
Schur expansion.

\section{The classical theorem at $\tau=1$}

The 2026 paper states the classical insertion theorem using the statistic
\[
  \di(w)=\#\{(i,j):i<j,\ w_i=w_j+1\}.
\]
For a fixed multiset \(S\), this is equivalent to the \(\tau=1\) version of
\(\dinv_\tau\), up to a constant shift.  Indeed, if
\(m_a\) is the multiplicity of the value \(a\) in \(S\), then every
rearrangement \(w\) of \(S\) has
\[
  \dinv_1(w)=\di(w)+\sum_a \binom{m_a}{2}.
\]
Thus grouping by \(\di\) or by \(\dinv_1\) carries the same information once
\(S\) is fixed.

The theorem-level result is the following Schur expansion.

\begin{theorem}[Classical Dyck symmetric functions]
Let \(S\) be a finite multiset of integers and let \(d\ge0\).  In the
classical step-one Dyck setting, the dual Dyck symmetric function has expansion
\[
  \DSstar(S,d;\mathbf x)
  =
  \sum_P s_{\lambda(P)}(\mathbf x),
\]
where \(P\) ranges over Dyck tableaux with entries \(S\) and
\(\di(\RR(P))=d\).  The affine Dyck symmetric function has the conjugate-shape
expansion
\[
  \DS(S,d;\mathbf x)
  =
  \sum_P s_{\lambda(P)'}(\mathbf x),
\]
with the same indexing set.
\end{theorem}

The dual formula is proved by an explicit insertion bijection.  A dual Dyck
factorization \(\mathcal F=(F_0,F_1,\ldots)\) is sent to a pair \((P,Q)\),
where \(P\) is a Dyck tableau and \(Q\) is a semistandard Young tableau of the
same shape.  The content of \(Q\) records the factor lengths:
\[
  m_i(Q)=|F_i|.
\]
Therefore the monomial \(x^{\mathcal F}\) becomes the ordinary SSYT monomial
attached to \(Q\).  Summing over all semistandard recording tableaux of fixed
shape gives the Schur function \(s_{\lambda(P)}(\mathbf x)\).

The affine formula is not proved by a separate nice affine insertion algorithm
in the paper.  It is derived from the dual formula using the standard
involution \(\omega\) on symmetric functions, together with
\(\omega(s_\lambda)=s_{\lambda'}\).  This asymmetry is important: the explicit
insertion algorithm currently belongs to the dual side.

\section{Why the insertion algorithm is not ordinary row insertion}

The classical dual insertion algorithm is close in spirit to dual RSK, but it
is not just ordinary row insertion with different notation.  Its distinctive
feature is that it sometimes inserts chunks rather than individual letters.
When the next input letter is exactly one more than a row entry, the algorithm
compares the maximal contiguous \(+2\)-chain in the row with the maximal
contiguous \(+2\)-chain at the front of the current input.  Depending on which
chain is shorter, an entire block is either swapped with a row block or passed
upward.

This chunking depends on the factorization, not merely on the underlying
concatenated word.  The paper's four-factor example is
\[
  F_0|F_1|F_2|F_3
  =
  (0,2,4)\,|\,(1,3)\,|\,(1,3,5)\,|\,(0,6).
\]
When \(F_1=(1,3)\) is inserted into the row \((0,2,4)\), the row has a longer
\(+2\)-chain than the input, so the input block passes through and the row is
unchanged.  Later, when \(F_2=(1,3,5)\) is inserted, the same prefix
\((1,3)\) belongs to a longer input chain; now the input chain has enough
length to displace the row chain.  This illustrates why two factorizations
with the same underlying word can lead to different insertion histories and,
in general, different insertion tableaux.  The factor boundaries are part of
the data.

This behavior is also what makes the algorithm interesting as a Dyck
deformation of dual RSK.  At \(\tau=0\), dual factors are simply strictly
increasing words, and the Schur expansion is governed by the classical dual
RSK correspondence.  At \(\tau=1\), strict increase is replaced by the dual
Dyck gap condition \(w_{i+1}\ge w_i+2\), and ordinary bumping must be replaced
by chunked row insertion.  Finding an analogous insertion procedure for
\(\tau>1\) remains open.

\section{The step-zero baseline}

The \(\tau=0\) specialization is not new combinatorics; it is the RSK/dual-RSK
baseline.  In this case \(d_0(a,b)=0\) for all pairs \((a,b)\), so the dinv
parameter has only the value \(0\) for each fixed multiset.  Dual factors
satisfy
\[
  w_{i+1}>w_i,
\]
so they are strictly increasing.  The dual Schur expansion is therefore the
usual dual-RSK Schur expansion for strict factor words.  The affine condition
becomes
\[
  w_{i+1}\le w_i,
\]
which is the weak/ordinary RSK counterpart after reversing the row convention.
This is why the \(\tau=1\) insertion theorem should be viewed as the first
nontrivial Dyck deformation of a respected classical algorithm rather than as
an isolated construction.

\section{The $r=s\tau+1$ conjecture}

The natural rational conjecture keeps the same tableau indexing set but
replaces the step-one inequalities and statistic by their step-\(\tau\)
versions.

\begin{conjecture}[Rational Dyck symmetric functions for \(r=s\tau+1\)]
Let \(\tau\ge0\), let \(S\) be a finite multiset of integers, and let
\(d\ge0\).  Then
\[
  \operatorname{DS}^{\ast,(\tau)}(S,d;\mathbf x)
  =
  \sum_P s_{\lambda(P)}(\mathbf x),
\]
and
\[
  \DS^{(\tau)}(S,d;\mathbf x)
  =
  \sum_P s_{\lambda(P)'}(\mathbf x),
\]
where \(P\) ranges over rational Dyck tableaux of step \(\tau\) with entries
exactly \(S\) and
\[
  \dinv_\tau(\RR(P))=d.
\]
\end{conjecture}

The conjecture is known at \(\tau=0\) by RSK and dual RSK, and at \(\tau=1\)
by the 2026 Dyck insertion theorem, after the constant shift between \(\di\)
and \(\dinv_1\) for fixed multisets.  For \(\tau>1\), the statement is
conjectural.  The present repository contains finite computational checks, not
a proof.

\section{Computational evidence and the role of \texttt{code.py}}

The file \texttt{code.py} is intended as a readable reference implementation
of the finite checks for this item.  It is not designed to be the fastest
possible implementation.  The larger recorded parameter boxes were checked
using more optimized code; the simplified file included here exposes the
mathematical comparison in a compact form.

For a fixed step \(\tau\), alphabet \(\{1,\ldots,A\}\), and maximum length
\(L\), the finite check proceeds as follows.  It enumerates word classes in the
bounded alphabet, groups them by multiset and \(\dinv_\tau\), enumerates the
valid affine and dual factorizations by composition type, and compares the
resulting coefficient counts with the Schur-side prediction obtained from
rational Dyck tableaux and semistandard Young tableaux.  Equivalently, for
each multiset and \(\dinv_\tau\) value, the check verifies that the
factorization counts agree with the Schur expansions predicted by the
conjecture.

The exhaustive parameter boxes recorded for the repository are
\[
  (\tau,A,L)=(2,10,10),\qquad (3,13,9),\qquad (4,16,8).
\]
In addition, sampled class checks were recorded for 100 classes in each of
\[
  (\tau,A,L)=(2,11,12),\qquad (3,14,11),\qquad (4,17,10).
\]
These computations support the conjecture for \(\tau>1\), but they do not
supply an insertion proof.  The open problem is to find a conceptual bijection
that explains the rational Schur expansions in the same way that row insertion
explains the classical dual case.

\section{Open directions}

Two insertion problems stand out.  First, one would like a genuine
step-\(\tau\) insertion algorithm for \(\tau>1\) that specializes to dual RSK
at \(\tau=0\) and to the chunked Dyck insertion algorithm at \(\tau=1\).  Such
an algorithm should explain why rational Dyck tableaux govern the Schur
expansion, rather than merely verifying the equality by enumeration.

Second, even in the classical \(\tau=1\) case, the known nice insertion
algorithm proves the dual formula directly.  The affine formula is obtained by
applying \(\omega\).  It would be interesting to have an explicit affine Dyck
insertion algorithm whose output directly produces the conjugate-shape
expansion.

\begin{thebibliography}{9}

\bibitem{Hawkes2026DyckSymmetric}
Graham Hawkes,
\emph{Dyck Symmetric Functions and Applications to \(q,t\)-Catalan
Polynomials},
arXiv:2605.13003, 2026.

\end{thebibliography}

\end{document}
